\section{QUBO Formulation}

\subsection{One-to-one objective with penalties}

The maximization in Eq.~\eqref{eq:objective} is represented as minimization. For the one-to-one case, the implemented energy is
\begin{equation}
Q_{1:1}(x)=
-\sum_{i,j}S_{ij}x_{ij}
+\lambda_A\sum_i\left(\sum_jx_{ij}-1\right)^2
+\lambda_R\sum_j\sum_{i<k}x_{ij}x_{kj}.
\label{eq:qubo_one_to_one}
\end{equation}
The first penalty enforces exactly one assignment in each agent row. The second penalizes resource collisions. In square instances, the combination yields permutation matrices at the feasible minimum.

\subsection{Exact-capacity QUBO}

For a capacity vector $\bm C$, the resource-side penalty becomes
\begin{equation}
Q_C(x)=
-\sum_{i,j}S_{ij}x_{ij}
+\lambda_A\sum_i\left(\sum_jx_{ij}-1\right)^2
+\lambda_R\sum_j\left(\sum_ix_{ij}-C_j\right)^2.
\label{eq:qubo_capacity}
\end{equation}
This exact squared penalty is the minimal generalization requested in the depth supplement \cite{depthsupplement}.

\subsection{Penalty scale and coefficient expansion}

The student-developed deterministic scale is retained:
\begin{equation}
\lambda_A=\lambda_R=nmS_{\max},
\qquad
S_{\max}=\max_{i,j}|S_{ij}|.
\label{eq:penalty_scale}
\end{equation}
For the three instances this gives Table~\ref{tab:penalties}.

\begin{table}[H]
\centering
\caption{Penalty weights and constant offsets.}
\label{tab:penalties}
\begin{tabular}{lrrrr}
\toprule
Instance & $N=nm$ & $\lambda_A$ & $\lambda_R$ & Constant offset\\
\midrule
$3\times3$ one-to-one & 9 & 8.55 & 8.55 & 25.65\\
$4\times4$ one-to-one & 16 & 14.72 & 14.72 & 58.88\\
$4\times2$, $\bm C=(2,2)$ & 8 & 7.20 & 7.20 & 86.40\\
\bottomrule
\end{tabular}
\end{table}

Writing a generic upper-triangular \QUBO{} as
\begin{equation}
Q(x)=c+\sum_qa_qx_q+\sum_{q<r}b_{qr}x_qx_r,
\label{eq:generic_qubo}
\end{equation}
the exact-capacity expansion yields
\begin{align}
c&=n\lambda_A+\lambda_R\sum_jC_j^2,\\
a_{ij}&=-S_{ij}-\lambda_A+\lambda_R(1-2C_j),\\
b_{(i,j),(i,\ell)}&=2\lambda_A,\qquad j<\ell,\\
b_{(i,j),(k,j)}&=2\lambda_R,\qquad i<k.
\end{align}
The coefficients are produced programmatically; no manually typed \QUBO{} matrix is used in the quantum stage.

\subsection{QUBO-to-Ising mapping}

Each binary variable is mapped to a Pauli-$Z$ operator through
\begin{equation}
x_q=\frac{1-Z_q}{2},
\qquad
x_qx_r=\frac{I-Z_q-Z_r+Z_qZ_r}{4}.
\label{eq:binary_ising_map}
\end{equation}
Thus,
\begin{equation}
H_Q=\alpha I+\sum_qh_qZ_q+\sum_{q<r}J_{qr}Z_qZ_r,
\label{eq:ising_hamiltonian}
\end{equation}
where
\begin{align}
\alpha&=c+\frac12\sum_qa_q+\frac14\sum_{q<r}b_{qr},\\
h_q&=-\frac12a_q-\frac14\sum_{r\neq q}b_{qr},\\
J_{qr}&=\frac14b_{qr}.
\end{align}
The notebook constructs the operator with explicit sparse Pauli terms and verifies basis-state equivalence between $Q(x)$ and $\bra{x}H_Q\ket{x}$. The validation is exhaustive for all 512 states of the $3\times3$ instance and all 256 states of the $4\times2$ instance; the $4\times4$ case is checked on deterministic reference states, optimum assignments, and control samples.

\subsection{Logical interaction topology}

The quadratic row and column couplings produce a dense rook-graph structure over the one-hot variables. Figure~\ref{fig:logical_topology} shows the logical graphs for the two square instances. This density is relevant because the IBM hardware coupling graph is sparse; routing therefore introduces SWAP-equivalent overhead even when the logical qubit count remains modest.

\begin{figure}[H]
\centering
\includegraphics[width=0.92\textwidth]{logical_qubo_topology.png}
\caption{Logical \QUBO{} interaction topology retained from the executed notebook. Row and column cliques create dense couplings before hardware mapping.}
\label{fig:logical_topology}
\end{figure}
